\documentclass[11pt,a4paper]{letter}
\usepackage[T1]{fontenc}
\usepackage[margin=2cm]{geometry}
\usepackage{lmodern}

\signature{Dmitriy Zhuravlev}
\address{Dmitriy Zhuravlev\\Taras Shevchenko National University of Kyiv\\Kyiv, Ukraine\\dzhuravlev@ukr.net}
\date{5 November 2026}

\begin{document}
\begin{letter}{The Editor-in-Chief\\EURASIP Journal on Image and Video Processing}
\opening{Dear Editor,}

I am pleased to submit the manuscript ``Learning-Free 3D Triangulation of Moving Objects from a Single Roadside Camera'' for consideration
as a Research Article in the EURASIP Journal on Image and Video Processing.

Monocular 3D perception for roadside cameras is today almost always a learned detector that outputs 3D boxes and
needs thousands of labelled boxes from the same kind of camera. The manuscript presents a learning-free 3D
triangulation of the objects that move past a fixed, calibrated, elevated camera. Its input is a mask and a 3D cuboid of each moving object per frame, from any source. The orientation of each object comes from its motion; in short windows
of the video its masked surface is triangulated, its planar rigid motion and the depth of every vertex solved
jointly, with the scale from the cuboid; the windows of one object are then solved together as one rigid body,
mirrored about its symmetry plane, and completed by the visual hull of its masks. Only symmetry is assumed, so the
same steps reconstruct cars, cyclists, motorcyclists and pedestrians. A classical way to obtain them from the video is described and evaluated; any other can replace it. On 237 cars of two public datasets, recorded at seven intersections, the models are
ahead of the labelled 3D boxes scored the same way against the cars' accumulated LiDAR on all four sets, with paired
bootstrap intervals clear of zero (by $+0.064$ on DAIR-V2X-C and $+0.026$ on V2X-Seq-SPD, also on the clips that
informed no design choice). The evaluation also shows that the raw accumulated LiDAR is a strongly biased reference at
long range, against which the labelled box scores only 0.35, and the manuscript reports its limits openly: slow cars
come out too long, and a walking pedestrian is triangulated only by its silhouettes.

I believe the work suits the journal because it is image and video processing in the classical sense: multi-view
geometry, dense optical flow, silhouette extraction and fusion of evidence over a whole video, applied to traffic
video and measured on public datasets with a reference that is itself examined critically. It offers readers an
interpretable alternative to learned detectors and a concrete account of where such geometry succeeds and fails.

The manuscript is original, has not been published and is not under consideration for publication elsewhere. My
earlier papers, cited in the manuscript, recovered 3D \emph{boxes} from single images with inverse perspective mapping
and segmented moving objects by optical-flow clustering --- methods that can supply the masks and cuboids the present
work takes as input; its 3D triangulation of the moving objects over the whole video and its LiDAR evaluation are
new. I am the sole author and approve the submission. I have no competing interests.

As the corresponding author is based in Ukraine, I would like to request a waiver of the article-processing charge
under the Springer Nature APC waiver policy.

Thank you for considering the manuscript.

\closing{Yours sincerely,}
\end{letter}
\end{document}
